% !TEX root = ../main.tex
\section{Discussion}
\label{sec:discussion}

\subsection{Hypotheses tested and refuted}

Hopper's inverted scaling (Section~\ref{sec:hopper}) absorbed the larger part
of our experimental effort. We record below the hypotheses that measurement
subsequently refuted. They are worth setting out rather than omitting, both
because each one narrows the space in which the remaining explanation must sit,
and because the order in which they failed is informative about where the
intuitions behind this method mislead.

\begin{description}
  \item[ADMM consensus quality.] At $N=8$ the agents' $\by_i$ are mutually
    near-orthogonal (mean pairwise cosine $\approx 0.05$), and their average
    retains only $\approx 0.57$ alignment with the true pooled natural-gradient
    direction against $0.75$--$0.83$ for a single agent. Forcing consensus with five ADMM rounds per policy iteration raised the
    cosine to $\approx 0.76$ and made performance \emph{worse}, and an exact
    pooled solve, which has zero consensus error by construction, landed within
    noise of one round. Consensus quality is not the bottleneck.
  \item[Trust-region overshoot.] Line~12's closed-form step is exact only at
    ADMM convergence, so with one round it could in principle overshoot. The
    measured overshoot factor was $0.68$--$1.04$ and the achieved KL
    $0.0054$--$0.0107$ against $\delta = 0.01$, so the step is mildly
    \emph{conservative} at both $N=1$ and $N=8$.
  \item[Architecture.] Table~4 mandates ReLU hidden layers and a Tanh-squashed
    policy mean, whereas no reference implementation squashes the mean. Testing
    all four combinations, the Table~4 configuration won, and removing the Tanh
    alone was the worst. The natural gradient largely absorbs the
    reparameterization, so the apparent gradient attenuation does not bite.
  \item[Exploration noise.] Three interventions left the fall rate at
    $9$--$12\%$ (Section~\ref{sec:variance}): releasing the log-std floor,
    releasing the ceiling, and annealing the noise to a level at which a fixed
    policy never falls.
\end{description}

The explanation that survives is the reward-structure argument of
Section~\ref{sec:hopper}, which is consistent with the known tendency of
MuJoCo survival bonuses to create local optima \citep{mania2018ars,
reda2020locomote}.

\subsection{Two observations we could not find in the literature}

We located no published source for either of the following, and present them
as our own measurements. First, a quantified failure rate and bimodal return
distribution for a trained stochastic MuJoCo policy. The literature states
that the episode ends when the robot falls, but not how often that happens to
a converged policy. Second, per-dimension divergence of a Gaussian policy's
$\log\sigma$: with the floor released, one action dimension collapsed to
$0.056$ while another drifted up to $0.311$, leaving the mean unchanged. TRPO
has no entropy term and the trust region bounds the KL of the joint
distribution, so a dimension whose actions barely affect the advantage
receives almost no gradient on its $\log\sigma$.

\subsection{Threats to validity}

Our study uses five seeds against the original's ten, so our error bars are
wider than they appear relative to a ten-seed estimate. All runs use a single
machine and a single software stack, so implementation-level differences from
the authors' setup cannot be separated from algorithmic ones. That concern has
precedent here: the same algorithm under identical hyperparameters has been
shown to differ by more than a factor of two on Hopper across implementations
\citep{henderson2018matters}. Finally, two hyperparameters ($\eta$ and the
interpretation of Table~4's learning-rate column,
Section~\ref{sec:implementation}) are our reading of an ambiguous source, and
a different reading could change the Hopper conclusion in particular.
